% Generated draft: requires mathematical and source review.
\documentclass[11pt]{article}
\usepackage[margin=25mm]{geometry}
\usepackage{fontspec}
\setmainfont{Latin Modern Roman}
\usepackage{amsmath,amssymb,mathtools,unicode-math}
\setmathfont{Latin Modern Math}
\usepackage{enumitem,needspace,xurl,hyperref}
\setlength{\parindent}{0pt}
\setlength{\parskip}{5pt}
\setlength{\emergencystretch}{3em}
\newcommand{\N}{\mathbb{N}}
\newcommand{\Z}{\mathbb{Z}}
\newcommand{\Q}{\mathbb{Q}}
\newcommand{\R}{\mathbb{R}}
\newcommand{\C}{\mathbb{C}}
\newcommand{\F}{\mathbb{F}}
\newcommand{\A}{\mathbb{A}}
\newcommand{\PP}{\mathbb{P}}
\newcommand{\E}{\mathbb{E}}
\newcommand{\eps}{\varepsilon}
\newcommand{\dd}{\,\mathrm{d}}
\begin{document}
\section*{UnsolvedMath ID: 30003256}
\textbf{Source problem code:} catalogue entry
\subsection{Definitions and mathematical statement}
The supplied record contains two intended questions, with unresolved signature and interpretation issues described below.

For the first question, let \(R\) be a Euclidean ring with Euclidean function \(\delta\). The intended preorder is \(a\preccurlyeq b\) when \(\delta(a)\leq\delta(b)\), wherever these values are defined. Let \(r,s\in R\) be multiplicatively independent noninvertible elements. For \(t\in\{r,s\}\), the source describes \(V_t(a)\) as the highest power of \(t\) dividing \(a\), and \(\mathbin{\cdot_t}\) as a multiplication relation restricted according to powers of \(t\). It does not supply complete definitions of these symbols.

The clean statement asks for a characterization of
\[
\operatorname{Def}_1(\mathcal R_r)\cap\operatorname{Def}_1(\mathcal R_s),
\qquad
\mathcal R_t=(R,+,-,0,\preccurlyeq,V_t,\mathbin{\cdot_t}),
\]
where \(\operatorname{Def}_1(\mathcal M)\) denotes the collection of subsets of the domain definable by first-order formulas with one free variable, under a parameter convention that the source does not specify. Thus the intended objects are subsets \(A\subseteq R\) admitting definitions in both structures, potentially by different formulas.

The original statement instead lists
\[
(R,+,-,0,\mathrm{multiplication},V_t,\mathbin{\cdot_t}).
\]
This differs materially from the clean signature. The supplied verification label favors the clean statement, but the record alone does not resolve whether ordinary multiplication or the preorder is intended. No identification of the two signatures is asserted here. The meaning of Euclidean ring, the domain and zero convention of \(\delta\), the precise meaning of multiplicative independence, the behavior of \(V_t\) at zero and at elements without a greatest dividing power, and the restriction defining \(\mathbin{\cdot_t}\) all require clarification before this becomes a fully specified characterization problem.

For the second question, with \(p\) prime as required by the notation \(\mathbb F_p\), set
\[
A_p=\left\{f\in\mathbb F_p[[X]]:\text{ there exists }0\ne P(T)\in\mathbb F_p(X)[T]\text{ such that }P(f)=0\right\},
\]
where evaluation takes place in \(\mathbb F_p((X))\). The supplied ambient structure is
\[
\mathcal F=\bigl(\mathbb F_p[[X]],+,0,V_X,\lambda_X,\prec,(\mathbin{\cdot_C})_{C\in\mathbb F_q[X]}\bigr),
\]
and \(\mathcal F_0\) is described as its substructure with domain \(A_p\). The intended question is whether the first-order theory
\[
\operatorname{Th}(\mathcal F_0)=\{\sigma:\sigma\text{ is a sentence in the intended language and }\mathcal F_0\models\sigma\}
\]
is decidable: does an algorithm halt on every sentence in an effective presentation of that language and correctly determine its truth in \(\mathcal F_0\)?

The source does not define \(V_X\), \(\lambda_X\), \(\prec\), or the symbols \(\mathbin{\cdot_C}\) completely, including their arities and boundary conventions. It also leaves \(q\) and its relationship to \(p\) unspecified. Consequently, the displayed notation does not yet determine a unique structure or decision problem. Calling \(A_p\) a substructure additionally requires closure under every intended function symbol; this is not established by the supplied record.
\subsection{Short English statement}
The first question asks which subsets of a Euclidean ring are definable in each of two expansions associated with multiplicatively independent nonunits. The second asks whether the first-order theory of an expansion of the algebraic power series over \(\mathbb F_p\) is decidable. Conflicting signatures and missing symbol definitions prevent a fully specified formulation. The supplied curation labels both questions open; that status is not independently verified here.
\subsection{Sources}
\begin{itemize}
\item Supplied local registry record: \texttt{data/unsolvedmath/problems.json}, permanent UnsolvedMath ID \(30003256\), exact source problem code \texttt{catalogue entry}, title \emph{Definability in Automata-Presentable Ring Expansions}. No individual proposer is named; the proposed year is recorded as 2016.
\item Supplied report citation: \emph{Definability and Decidability Problems in Number Theory} (2017), Oberwolfach Reports, pp. 2793--2866, DOI 10.4171/owr/2016/49. \url{https://doi.org/10.4171/owr/2016/49}. The record attributes the problem to workshop open-problem discussions documented in this report. The 2016 proposed year and the supplied 2017 report designation are retained without reconciliation.
\item Supplied literature source: Point--Rigo--Waxweiler, \emph{Polynomial rings}. \url{https://doi.org/10.1080/00927872.2015.1028236}.
\item Supplied literature source: \emph{Oberwolfach Reports 2017}. \url{https://ems.press/content/serial-article-files/46652}.
\item Literature assessment checked at \texttt{2026-08-21}: the supplied curation labels the common-definability characterization and decidability question open and states that known polynomial-ring expansions do not settle them. This is supplied dated curation, not an independently verified literature assessment. Its own qualification requests independent verification of time-sensitive status and cited claims.
\item The registry labels the statement \texttt{corrected\_verified} and reports comparison against the official Oberwolfach report. These are supplied verification claims; no source visit or independent verification is asserted in this draft.
\end{itemize}
\end{document}
